% problem-000030 9 二吡喃酮热裂解合成
\section*{9. 二吡喃酮热裂解合成}
2-吡喃酮(N)的衍⽣物⼴泛存在于⾃然界，有强⼼、抑菌等作⽤。化合物L真空热裂解，经中间体M转化为N。M是具有⼀定稳定性和强反应活性的⾮环状分⼦。

（图：）

利⽤该反应原理，起始原料O通过系列反应，成功合成2-吡喃酮衍⽣物 $\mathbf { T } _ { \circ } \ \mathbf { 0 } \to \mathbf { R }$ 的反应具有很好的⽴体选择性。

（图：）

$
\underset {\left(\mathrm{C} _ {1 7} \mathrm{H} _ {1 6} \mathrm{O} _ {2}\right)} {\mathbf {R}} \xrightarrow [ \mathrm{DCM/MeOH,rt} ]{\mathrm{H} _ {2} \mathrm{O} _ {2} , \mathrm{NaOH}} \quad \mathbf {S} \xrightarrow [ 6 5 0 ^ {\circ} \mathrm{C} ]{\text {FVT}} \underset {\left(\mathrm{C} _ {1 2} \mathrm{H} _ {1 0} \mathrm{O} _ {3}\right)} {\mathbf {T}} \quad \text {备注：DCM = CH_{2} Cl_{2}}
$

\subsection*{9-1}
请写出M及Q、R、S、T的结构简式（如为⽴体选择性反应产物，需表明⽴体化学结构）。

\subsection*{9-2}
判断化合物N有⽆芳⾹性，并简述理由。

\subsection*{9-3}
请⽐较化合物N及1,3-环⼰⼆烯作为双烯体发⽣Diels‒Alder反应的相对活性⼤⼩，并简述理由。

\subsection*{9-4}
请指出在Q→R反应中，PhOH的作⽤。
